\subsection{one-hot法}
\label{subsec:lab4-one-hot}

\paragraph{基本原理}

一个最直接的想法就是，我们直接就给所有点一个变量，集合的属于其实就是0,1关系嘛，0就是不属于，1就是属于

所以说我们给每个点一个0,1变量,记作$\lambda_{(a,b)}$，如果点$(a,b)$属于点集,那么$\lambda_{(a,b)}=1$，否则$\lambda_{(a,b)}=0$

那你求点集里的点，其实就是\(\lambda_{(a,b)}=1\),然后其他都是0

所以约束条件就出来了，\[
\sum_{(a,b)\in V}\lambda_{(a,b)}=1,
\]

\paragraph{数学建模}
\begin{equation}
\begin{aligned}
\min\quad &0\\
\text{s.t.}\quad
&\sum_{(a,b)\in V}\lambda_{(a,b)}=1\\
&x=\sum_{(a,b)\in V}\lambda_{(a,b)}\operatorname{bits}(a,b)\\
&\lambda_{(a,b)}\in\{0,1\},\forall(a,b)\in V
\end{aligned}
\label{eq:lab4-one-hot-model}
\end{equation}


\paragraph{正确性分析}

对于模型的任意可行解，由 $\lambda_{(a,b)}\in\{0,1\}$ 以及
\[
\sum_{(a,b)\in V}\lambda_{(a,b)}=1
\]
可知，存在唯一的 $(a^*,b^*)\in V$，使得
$\lambda_{(a^*,b^*)}=1$，其余选择变量均为 $0$

将其代入坐标约束
\[
x_j=\sum_{(a,b)\in V}
\operatorname{bits}(a,b)_j\lambda_{(a,b)}
\]
可得
\[
x_j=\operatorname{bits}(a^*,b^*)_j,
\quad j=0,1,\ldots,7.
\]
因此，$x=\operatorname{bits}(a^*,b^*)$，即 $x$ 对应于 $V$ 中的点，不会有不可达的点

反过来，任取 $(a,b)\in V$，令
\[
x=\operatorname{bits}(a,b),\qquad
\lambda_{(a,b)}=1,
\]
其余选择变量均为 $0$，则显然满足模型中的全部约束

\paragraph{用于差分与线性传播}

差分建模时取$V=V_{\mathrm D}$，每个选择变量对应一个DDT非零位置；线性建模时取$V=V_{\mathrm L}$，每个选择变量对应一个LAT非零位置。也可以直接用$(a,b)$给选择变量编号，记作$\lambda_{(a,b)}$，从而与原表中的位置一一对应。两类模型分别建立各自的选择变量，坐标约束形式保持一致；全零点的选择变量也需要保留，以表示未活跃S盒。表项的具体数值不作为这里的优化目标，枚举结果应覆盖全部可达点。
